
\documentclass[addpoints]{exam}
\usepackage{amsmath}
\usepackage{amssymb}
\usepackage{color}
\usepackage{siunitx}

% \printanswers

\newcommand{\event}{Trigonometric Ratios}
\newcommand{\tournament}{}
\def\type{0}

\setlength\fillinlinelength{\textwidth}
\CorrectChoiceEmphasis{\color{red}\bfseries}
\SolutionEmphasis{\color{red}\bfseries}

\setlength\answerclearance{2pt}
\pagestyle{headandfoot}
\runningheadrule
\runningheader{\event}{\tournament}{\ifnum\type=1 Class Set Test \else Full Name:\hspace{1cm} \fi}
\runningfootrule
\runningfooter{}{Page \thepage\ of \numpages}{}

\begin{document}
\begin{center}
    {\huge Grade 11 Functions \\ \event
    \ifprintanswers \space Key
    \else \space \ifnum\type<2 Test \fi \ifnum\type=1 \textbf{(CLASS SET)} \fi
          \ifnum\type=2 Answer Sheet \fi
    \fi \\ \vspace{3mm}\tournament\vspace{1mm}}
\end{center}

\vspace{.25\textheight}

\begin{center}
    First Name: \ifprintanswers
    \underline{\hspace{3.5cm}}\textbf{KEY}\underline{\hspace{5.5cm}}\\\vspace{5mm}
    \else \underline{\hspace{10cm}}\\ \vspace{5mm} \fi
    Last Name: \underline{\hspace{10cm}}\\ \vspace{5mm}
\end{center}

\noindent\textbf{\underline{Directions:}}
\begin{itemize}
    \ifnum\type=1 \item \textbf{This test is a class set.} Do not write on this paper. \fi
    \item Show all work for full marks. Give exact values where possible.
    Diagrams required for ambiguous case.
    \item The test is designed to be completed in 75 minutes.
\end{itemize}
\begin{center}
\textit{For grading use only}\\
\multirowgradetable{1}[pages]
\end{center}

\newpage
\begin{questions}

\section*{Part A --- Multiple Choice (15 marks)}
\vspace{5mm}

\question[1] The exact value of $\sin\ang{45}$ is:
\begin{choices}
    \choice $\dfrac{1}{2}$
    \choice $\dfrac{\sqrt{3}}{2}$
    \correctchoice $\dfrac{\sqrt{2}}{2}$
    \choice $1$
\end{choices}

\question[1] The exact value of $\cos\ang{30}$ is:
\begin{choices}
    \choice $\dfrac{1}{2}$
    \correctchoice $\dfrac{\sqrt{3}}{2}$
    \choice $\dfrac{\sqrt{2}}{2}$
    \choice $\sqrt{3}$
\end{choices}

\question[1] In which quadrant is $\theta = \ang{250}$ located?
\begin{choices}
    \choice Quadrant I
    \choice Quadrant II
    \correctchoice Quadrant III
    \choice Quadrant IV
\end{choices}

\question[1] If $\sin\theta < 0$ and $\cos\theta > 0$, then $\theta$ is in:
\begin{choices}
    \choice Quadrant I
    \choice Quadrant II
    \choice Quadrant III
    \correctchoice Quadrant IV
\end{choices}

\question[1] The reference angle for $\theta = \ang{210}$ is:
\begin{choices}
    \choice \ang{70}
    \choice \ang{60}
    \correctchoice \ang{30}
    \choice \ang{150}
\end{choices}

\question[1] $\cos\ang{120}$ equals:
\begin{choices}
    \choice $\dfrac{\sqrt{3}}{2}$
    \correctchoice $-\dfrac{1}{2}$
    \choice $\dfrac{1}{2}$
    \choice $-\dfrac{\sqrt{3}}{2}$
\end{choices}

\question[1] $\sin^2\theta + \cos^2\theta =$
\begin{choices}
    \choice $0$
    \choice $\sin\theta\cos\theta$
    \correctchoice $1$
    \choice $2$
\end{choices}

\question[1] If $\cos\theta = -\dfrac{3}{5}$ and $\theta$ is in Quadrant III, then
$\sin\theta =$
\begin{choices}
    \choice $\dfrac{4}{5}$
    \correctchoice $-\dfrac{4}{5}$
    \choice $\dfrac{3}{4}$
    \choice $-\dfrac{3}{4}$
\end{choices}

\question[1] The Sine Law states:
\begin{choices}
    \choice $c^2 = a^2 + b^2 - 2ab\cos C$
    \correctchoice $\dfrac{a}{\sin A} = \dfrac{b}{\sin B} = \dfrac{c}{\sin C}$
    \choice $\sin A = \dfrac{a}{c}$ only for right triangles
    \choice $a^2 = b^2 + c^2$
\end{choices}

\question[1] The Cosine Law $c^2 = a^2 + b^2 - 2ab\cos C$ reduces to the Pythagorean
theorem when:
\begin{choices}
    \choice $C = \ang{60}$
    \choice $a = b$
    \correctchoice $C = \ang{90}$
    \choice $\cos C = 1$
\end{choices}

\question[1] In $\triangle ABC$, $A = \ang{30}$, $a = 5$, $b = 8$. The ambiguous case arises
because:
\begin{choices}
    \choice $a > b$
    \correctchoice $b\sin A < a < b$, so two triangles may exist
    \choice $A > \ang{90}$
    \choice $a < b\sin A$, so no triangle exists
\end{choices}

\question[1] $\tan\ang{135}$ equals:
\begin{choices}
    \choice $1$
    \correctchoice $-1$
    \choice $\sqrt{3}$
    \choice $-\sqrt{3}$
\end{choices}

\question[1] In a triangle with sides $a = 7$, $b = 9$, $c = 5$, which formula
should be used to find angle $B$?
\begin{choices}
    \choice Sine Law — two sides and the included angle are known
    \correctchoice Cosine Law — all three sides are known
    \choice Pythagorean theorem
    \choice Sine Law — two angles are known
\end{choices}

\question[1] Two angles in $[0^\circ, 360^\circ]$ have $\sin\theta = \dfrac{\sqrt{3}}{2}$.
They are:
\begin{choices}
    \choice \ang{30} and \ang{150}
    \correctchoice \ang{60} and \ang{120}
    \choice \ang{30} and \ang{330}
    \choice \ang{60} and \ang{300}
\end{choices}

\question[1] $\sec\theta$ is defined as:
\begin{choices}
    \choice $\dfrac{1}{\sin\theta}$
    \correctchoice $\dfrac{1}{\cos\theta}$
    \choice $\dfrac{\cos\theta}{\sin\theta}$
    \choice $\dfrac{\sin\theta}{\cos\theta}$
\end{choices}


\section*{Part B --- Short Answer (33 marks)}

\question Angles in standard position and exact values.
\begin{parts}
    \part[3] For $\theta = \ang{315}$, find $\sin\theta$, $\cos\theta$, and $\tan\theta$
    using the reference angle. State the quadrant.
    \part[3] If $\sin\theta = \dfrac{5}{13}$ and $\theta$ is in Quadrant II, find
    $\cos\theta$ and $\tan\theta$ exactly.
    \part[2] Find all angles $\theta \in [0^\circ, 360^\circ]$ where
    $\cos\theta = -\dfrac{\sqrt{2}}{2}$.
\end{parts}
\begin{solution}[10cm]
\textbf{(a)} $\theta = \ang{315}$ is in Quadrant IV (reference angle $\ang{45}$).
\[
\sin\ang{315} = -\frac{\sqrt{2}}{2};\quad \cos\ang{315} = \frac{\sqrt{2}}{2};\quad
\tan\ang{315} = -1
\]

\textbf{(b)} $\sin\theta = \dfrac{5}{13}$, Quadrant II.
Using $\sin^2\theta + \cos^2\theta = 1$:
$\cos^2\theta = 1 - \dfrac{25}{169} = \dfrac{144}{169}$, so $\cos\theta = -\dfrac{12}{13}$
(negative in QII).
\[
\tan\theta = \frac{\sin\theta}{\cos\theta} = \frac{5/13}{-12/13} = -\frac{5}{12}
\]

\textbf{(c)} Reference angle: $\cos^{-1}\!\left(\dfrac{\sqrt{2}}{2}\right) = \ang{45}$.
Cosine is negative in QII and QIII.
$\theta = \ang{135}$ or $\theta = \ang{225}$
\end{solution}

\question Apply the Sine Law to find the unknown sides/angles. Round to one decimal place.
\begin{parts}
    \part[3] In $\triangle ABC$: $A = \ang{42}$, $B = \ang{75}$, $a = 15$ cm.
    Find side $b$.
    \part[4] \textbf{Ambiguous case:} In $\triangle ABC$: $A = \ang{35}$, $a = 9$ cm,
    $b = 14$ cm. Determine whether 0, 1, or 2 triangles exist. Solve all possible
    triangles.
\end{parts}
\begin{solution}[10cm]
\textbf{(a)} $\dfrac{b}{\sin B} = \dfrac{a}{\sin A}$:\quad
$b = \dfrac{a\sin B}{\sin A} = \dfrac{15\sin\ang{75}}{\sin\ang{42}}
= \dfrac{15 \times 0.9659}{0.6691} = \mathbf{\SI{21.7}{cm}}$

\textbf{(b)} Check: $b\sin A = 14\sin\ang{35} = 14(0.5736) = 8.03$.
Since $b\sin A < a < b$ ($8.03 < 9 < 14$), \textbf{two triangles exist}.

$\dfrac{\sin B}{b} = \dfrac{\sin A}{a}$:\quad
$\sin B = \dfrac{b\sin A}{a} = \dfrac{14\sin\ang{35}}{9} = \dfrac{8.030}{9} = 0.8923$

$B_1 = \ang{63.2}$, so $C_1 = 180 - 35 - 63.2 = \ang{81.8}$.
$c_1 = \dfrac{9\sin\ang{81.8}}{\sin\ang{35}} = \dfrac{9(0.9898)}{0.5736} \approx \mathbf{\SI{15.5}{cm}}$

$B_2 = 180 - 63.2 = \ang{116.8}$, so $C_2 = 180 - 35 - 116.8 = \ang{28.2}$.
$c_2 = \dfrac{9\sin\ang{28.2}}{\sin\ang{35}} \approx \dfrac{9(0.4726)}{0.5736}
\approx \mathbf{\SI{7.42}{cm}}$
\end{solution}

\question Apply the Cosine Law.
\begin{parts}
    \part[3] In $\triangle PQR$: $p = 8$, $q = 11$, $R = \ang{62}$.
    Find side $r$.
    \part[3] In $\triangle XYZ$: $x = 5$, $y = 7$, $z = 9$. Find angle $Z$.
\end{parts}
\begin{solution}[9cm]
\textbf{(a)} $r^2 = p^2 + q^2 - 2pq\cos R = 64 + 121 - 2(8)(11)\cos\ang{62}$
$= 185 - 176(0.4695) = 185 - 82.63 = 102.37$
\[
r = \sqrt{102.37} \approx \mathbf{\SI{10.1}{cm}}
\]

\textbf{(b)} $\cos Z = \dfrac{x^2 + y^2 - z^2}{2xy} = \dfrac{25 + 49 - 81}{2(5)(7)}
= \dfrac{-7}{70} = -0.1$
\[
Z = \cos^{-1}(-0.1) \approx \mathbf{\ang{95.7}}
\]
\end{solution}

\question A surveyor needs to find the width of a river. From point $A$ on one bank,
she measures \ang{65} to a tree $C$ on the opposite bank. She walks 80 m along the
bank to point $B$ and measures \ang{52} to the same tree. The angle at $C$ is
therefore $\angle ACB = 180^\circ - 65^\circ - 52^\circ = 63^\circ$.
\begin{parts}
    \part[3] Find the distance $AC$ using the Sine Law.
    \part[2] Find the perpendicular width of the river (the distance from $C$
    perpendicular to $AB$).
\end{parts}
\begin{solution}[8cm]
\textbf{(a)} In $\triangle ABC$: $AB = 80$ m, $\angle A = \ang{65}$,
$\angle B = \ang{52}$, $\angle C = \ang{63}$.
\[
\frac{AC}{\sin B} = \frac{AB}{\sin C} \implies
AC = \frac{80\sin\ang{52}}{\sin\ang{63}} = \frac{80(0.7880)}{0.8910}
\approx \mathbf{\SI{70.7}{m}}
\]

\textbf{(b)} Width $= AC \cdot \sin(\angle A) = 70.7 \times \sin\ang{65}
= 70.7 \times 0.9063 \approx \mathbf{\SI{64.1}{m}}$
\end{solution}

\end{questions}
\end{document}
